\section{Series de Fourier}
\subsection{Introducción}

\subsubsection{Contexto Histórico}
\begin{itemize}
    \item $1750$: D. Bernoulli propone $f(x) = \sum_{n=1}^{\infty} a_n \sin(nx)$ para la ecuación de onda. Intervienen Euler, D'Alembert, Lagrange.
    \item $\simeq 1800 - 1807$: Aparece la representación $f(x) = \frac{a_0}{2} + \sum_{n=1}^{\infty} (a_n \cos(nx) + b_n \sin(nx))$.
    \item $1822$: ¿$a_n, b_n$?
\end{itemize}
Integrando en un periodo para hallar $a_0$:
\begin{align*}
    \int_{-\pi}^{\pi} f(x) dx &= \int_{-\pi}^{\pi} \left[ \frac{a_0}{2} + \sum_{n=1}^{\infty} (a_n \cos(nx) + b_n \sin(nx)) \right] dx \\
    &= \int_{-\pi}^{\pi} \frac{a_0}{2} dx + 0 = \pi a_0
\end{align*}
De manera análoga, multiplicando por $\cos(mx)$ e integrando:
\begin{align*}
    \int_{-\pi}^{\pi} f(x) \cos(mx) dx &= \int_{-\pi}^{\pi} \left[ \frac{a_0}{2} + \sum_{n=1}^{\infty} (a_n \cos(nx) + b_n \sin(nx)) \right] \cos(mx) dx \\
    &= \int_{-\pi}^{\pi} \frac{a_0}{2} \cos(mx) dx \\
    &\quad + \sum_{n=1}^{\infty} \left( a_n \int_{-\pi}^{\pi} \cos(nx) \cos(mx) dx + b_n \int_{-\pi}^{\pi} \sin(nx) \cos(mx) dx \right) \\
    &= 0 + a_m \int_{-\pi}^{\pi} \cos^2(mx) dx + 0 \\
    &= a_m \pi \implies a_m = \frac{1}{\pi} \int_{-\pi}^{\pi} f(x) \cos(mx) dx
\end{align*}

Se puede desarrollar de forma análoga para $b_n$
\subsection{Funciones Periódicas}

\begin{definicion}{Función Periódica}{func_per}
$f: \mathbb{R} \to \mathbb{C}$ es periódica si $\exists T > 0$ tal que $f(x) = f(x+T)$ para todo $x \in \mathbb{R}$. En este caso, $T$ es el Periodo de $f$ ($f$ es $T$-periódica).
\end{definicion}

\begin{nota}{Construcción de funciones periódicas}{constr_per}
Para construir una función $T$-periódica basta con definir $\tilde{f}: [0, T[ \to \mathbb{C}$ y extenderla:
$$ f(x) = \tilde{f}(x - kT) \quad \text{si } kT \le x < (k+1)T, \; k \in \mathbb{Z} $$
También podemos considerar $\tilde{f}: [-T/2, T/2[ \to \mathbb{C}$ y \"copiar\".

\textbf{Ejemplos:}
\begin{enumerate}
    \item $f(x) = 1$ para $0 \le x < 1$, y $f(x) = -1$ para $-1 \le x < 0$. ($2$-periódica).
    \item $f(x) = x$ para $0 \le x < 1$. ($1$-periódica).
    \item $f(x) = x$ para $-1/2 \le x < 1/2$. ($1$-periódica).
\end{enumerate}
\end{nota}

\begin{observacion}{Cambio de Periodo}{cambio_per}
Pregunta: ¿Puede descomponerse cualquier función periódica en términos de funciones periódicas elementales?

Nota: Supongamos que $f: \mathbb{R} \to \mathbb{C}$ es $T$-periódica (la pienso como $f: [-T/2, T/2[ \to \mathbb{C}$ y extendida periódica). Tomo $L > 0$ y defino $g: [-L/2, L/2[ \to \mathbb{C}$ como $g(x) = f\left(\frac{T}{L}x\right)$.
$$ -\frac{L}{2} \le x < \frac{L}{2} \implies -\frac{T}{2} \le \frac{T}{L}x < \frac{T}{2} $$
¡El periodo no es relevante! A partir de ahora todas las funciones serán $2\pi$-periódicas:
$f: [-\pi, \pi[ \to \mathbb{C}$ o $f: [0, 2\pi[ \to \mathbb{C}$.
\end{observacion}

\subsection{Coeficientes y Series de Fourier}

\begin{definicion}{Coeficientes de Fourier}{coef_fourier}
Dada $f: \mathbb{R} \to \mathbb{C}$ $2\pi$-periódica, se definen sus coeficientes de Fourier (si las integrales existen):
\begin{align*}
    a_n &= \frac{1}{\pi} \int_{-\pi}^{\pi} f(x) \cos(nx) dx \quad n = 0, 1, 2, \dots \\
    b_n &= \frac{1}{\pi} \int_{-\pi}^{\pi} f(x) \sin(nx) dx \quad n = 1, 2, 3, \dots
\end{align*}
\end{definicion}

\begin{nota}{Integración en un periodo}{int_per}
Si $g: \mathbb{R} \to \mathbb{C}$ es $T$-periódica, entonces $\int_{-T/2}^{T/2} g(x) dx = \int_0^T g(x) dx$ (Para la demostración basta con hacer cambio de variable). En la definición de los coeficientes de Fourier se puede tomar $\int_0^{2\pi}$ o $\int_I$, con $I$ de longitud $2\pi$.
\end{nota}

\begin{definicion}{Serie de Fourier}{serie_fourier}
Si $f: \mathbb{R} \to \mathbb{C}$ es $2\pi$-periódica, se define su serie de Fourier (serie formal):
$$ f \sim \frac{a_0}{2} + \sum_{n=1}^{\infty} \left( a_n \cos(nx) + b_n \sin(nx) \right) $$
\end{definicion}

\subsubsection{Convergencia}

\begin{nota}{Tipos de Convergencia}{conv}
En un espacio normado $(V, ||\cdot||) \implies (v_n)_n$ converge a $v$ si $||v_n - v|| \xrightarrow{n \to \infty} 0$.

Para funciones $g_n: \Omega \to \mathbb{C}$:
\begin{itemize}
    \item \textbf{Conv puntual:} $|g_n(x) - g(x)| \to 0 \quad \forall x \in \Omega$
    \item \textbf{Conv unif:} $\sup_{x \in \Omega} |g_n(x) - g(x)| \to 0 \implies$ Intercambiar (límites, integrales, derivadas).
\end{itemize}
\end{nota}
\subsection{Forma compleja de la serie de Fourier}

A partir de las identidades de Euler:
\begin{equation*}
    \cos(nt) = \frac{e^{int} + e^{-int}}{2}, \quad \sin(nt) = \frac{e^{int} - e^{-int}}{2i} = -i \frac{e^{int} - e^{-int}}{2}
\end{equation*}

Entonces la serie se desarrolla como:
\begin{align*}
    \frac{a_0}{2} + \sum_{n=1}^{\infty} (a_n \cos(nt) + b_n \sin(nt)) &= \frac{a_0}{2} + \sum_{n=1}^{\infty} \left( \frac{a_n}{2} e^{int} + \frac{a_n}{2} e^{-int} - \frac{b_n i}{2} e^{int} + \frac{b_n i}{2} e^{-int} \right) \\
    &= \frac{a_0}{2} + \sum_{n=1}^{\infty} \frac{a_n - b_n i}{2} e^{int} + \sum_{n=1}^{\infty} \frac{a_n + b_n i}{2} e^{-int}
\end{align*}

Para el último sumatorio, aplicando el cambio de variable $-n=k$, se obtiene $\sum_{-\infty}^{-1} \frac{a_{-k} + b_{-k}i}{2} e^{ikt}$.

\begin{nota}{Notación}{notacion_coeficientes}
Se definen los coeficientes complejos de la siguiente manera:
\begin{itemize}
    \item Cuando $n = 0$: $\hat{f}(0) = \frac{a_0}{2}$
    \item Cuando $n \ge 1$: $\hat{f}(n) = \frac{a_n - i b_n}{2}$
    \item Cuando $n \le -1$: $\hat{f}(n) = \frac{a_{-n} + i b_{-n}}{2}$
\end{itemize}
\end{nota}

Entonces, la fórmula compleja de la serie de Fourier queda expresada como:
\begin{equation*}
    f(t) \sim \sum_{n=-\infty}^{\infty} \hat{f}(n) e^{int}
\end{equation*}

\subsubsection{Coeficientes de la forma compleja}

Desarrollando la expresión para cada rango de $n$:
\begin{itemize}
    \item Para $n = 0$:
    \begin{equation*}
        \hat{f}(0) = \frac{a_0}{2} = \frac{1}{2\pi} \int_{-\pi}^{\pi} f(t) dt
    \end{equation*}
    \item Para $n \ge 1$:
    \begin{align*}
        \hat{f}(n) &= \frac{1}{2}(a_n - ib_n) = \frac{1}{2\pi} \int_{-\pi}^{\pi} f(t) (\cos(nt) - i \sin(nt)) dt \\
        &= \frac{1}{2\pi} \int_{-\pi}^{\pi} f(t) (\overline{\cos(nt) + i \sin(nt)}) dt \\
        &= \frac{1}{2\pi} \int_{-\pi}^{\pi} f(t) \overline{e^{int}} dt = \frac{1}{2\pi} \int_{-\pi}^{\pi} f(t) e^{-int} dt
    \end{align*}
    \item Para $n \le -1$:
    \begin{equation*}
        \hat{f}(n) = \frac{1}{2}(a_{-n} + i b_{-n}) = \frac{1}{2\pi} \int_{-\pi}^{\pi} f(t) (\cos(nt) + i \sin(nt)) dt = \frac{1}{2\pi} \int_{-\pi}^{\pi} f(t) e^{-int} dt
    \end{equation*}
\end{itemize}

Los coeficientes complejos quedan de la forma general unificada:
\begin{equation*}
    \hat{f}(n) = \frac{1}{2\pi} \int_{-\pi}^{\pi} f(t) e^{-int} dt, \quad \forall n \in \mathbb{Z}
\end{equation*}

\begin{observacion}{Recuperación de coeficientes reales}{recuperacion_coeficientes}
Si conocemos la sucesión $(\hat{f}(n))_{n \in \mathbb{Z}}$, ¿cómo encontramos $a_n$ y $b_n$?  

Sabemos que $a_0 = 2\hat{f}(0)$. Despejando en las ecuaciones para $n \ge 1$, obtenemos:
\begin{align*}
    a_n &= \hat{f}(n) + \hat{f}(-n) \\
    b_n &= i(\hat{f}(n) - \hat{f}(-n))
\end{align*}
\end{observacion}

\subsection{El grupo topológico $\mathbb{T}$ y el espacio $C(\mathbb{T})$}

\begin{definicion}{El toro $\mathbb{T}$}{toro}
Se define el toro como $\mathbb{T} = \{z \in \mathbb{C} : |z|=1\} = \partial \mathbb{D}$.
\end{definicion}

\begin{observacion}{Propiedades de $\mathbb{T}$}{prop_toro}
El toro unidimensional $\mathbb{T} = \{z \in \mathbb{C} : |z|=1\} = \partial\mathbb{D}$ cumple las siguientes propiedades fundamentales:
\begin{enumerate}
    \item Como espacio topológico, es un subconjunto cerrado y acotado dentro de $\mathbb{C}$. Por tanto, es un espacio compacto.
    \item Tiene estructura de \textbf{grupo abeliano} bajo la multiplicación de números complejos. Es decir, para cualesquiera $z_1, z_2 \in \mathbb{T}$, se cumple que $|z_1 \cdot z_2| = |z_1| \cdot |z_2| = 1 \cdot 1 = 1$, por lo que $z_1 \cdot z_2 \in \mathbb{T}$, además $z_1^{-1} \in \mathbb{T}$ y el elemento neutro es $1 \in \mathbb{T}$.
\end{enumerate}
\end{observacion}

\begin{nota}{El homomorfismo $\Phi$ e isomorfismo con $\mathbb{R} / 2\pi\mathbb{Z}$}{homomorfismo_toro}
Considérese la aplicación $\Phi: \mathbb{R} \to \mathbb{T}$ definida como $\Phi(t) = e^{it}$. Esta función cumple lo siguiente:
\begin{enumerate}
    \item \textbf{Es un homomorfismo de grupos:} Para cualesquiera $t_1, t_2 \in \mathbb{R}$, la exponencial compleja verifica que:
    \begin{equation*}
        \Phi(t_1 + t_2) = e^{i(t_1+t_2)} = e^{it_1}e^{it_2} = \Phi(t_1)\Phi(t_2)
    \end{equation*}
    \item \textbf{Es sobreyectivo:} Para todo $z \in \mathbb{T}$, existe al menos un $t \in \mathbb{R}$ tal que $z = e^{it}$ (basta con tomar $t = \arg(z)$).
    \item \textbf{Núcleo de la aplicación:} El núcleo (o \textit{ker}) está formado por todos los valores de $t$ que se aplican en el elemento neutro de $\mathbb{T}$ (que es el $1$):
    \begin{equation*}
        \ker \Phi = \{t \in \mathbb{R} : e^{it} = 1\} = \{2\pi k : k \in \mathbb{Z}\} = 2\pi\mathbb{Z}
    \end{equation*}
\end{enumerate}
Por el Primer Teorema de Isomorfía de grupos, al ser $\Phi$ un homomorfismo exhaustivo, el dominio cocientado por el núcleo es isomorfo a la imagen. Entonces:
\begin{equation*}
    \mathbb{T} \cong \mathbb{R} / 2\pi\mathbb{Z}
\end{equation*}
Como consecuencia directa de este isomorfismo, todo elemento del toro puede representarse de forma única en un periodo; es decir, para todo $z \in \mathbb{T}$, existe un único $t \in [-\pi, \pi[$ tal que $z = e^{it}$.
\end{nota}

\begin{nota}{Funciones en el Toro}{func_toro}
Supongamos que tenemos $f:\mathbb{T} \to \mathbb{C}$ y definimos $g:\mathbb{R} \to \mathbb{C}$ como $g(t) = f(e^{it})$. La función $g$ es $2\pi$-periódica.

Inversamente, si $g:\mathbb{R} \to \mathbb{C}$ es $2\pi$-periódica, se puede definir $f(e^{it}) = g(t)$ con $f:\mathbb{T} \to \mathbb{C}$.

Además, $f$ es continua si y solo si $g$ es continua. De igual forma, $f$ es medible si y solo si $g$ es medible.

A partir de aquí, utilizaremos la notación $f(z) = f(e^{it}) = f(t)$ indistintamente para $z \in \mathbb{T}$, $t \in [-\pi,\pi[$.
\end{nota}

\begin{definicion}{El espacio $C(\mathbb{T})$}{espacio_ct}
Se define el espacio de funciones continuas sobre el toro como $C(\mathbb{T}) = \{f:\mathbb{T} \to \mathbb{C} \mid f \text{ continua}\}$.

Esto es equivalente a $\{f:[-\pi,\pi] \to \mathbb{C} \mid f(\pi)=f(-\pi) \text{ y continua}\}$ y a $\{f:\mathbb{R} \to \mathbb{C} \mid 2\pi\text{-periódica y continua}\}$.

Además, dotado de la norma del supremo $\|f\|_\infty = \sup_{z \in \mathbb{T}} |f(z)| = \sup_{t \in [-\pi,\pi[} |f(t)|$, el espacio $(C(\mathbb{T}), \|\cdot\|_\infty)$ es un espacio vectorial normado.
\end{definicion}

\subsection{Polinomios Trigonométricos y Densidad}

\begin{definicion}{Polinomio trigonométrico}{pol_trig}
Un polinomio trigonométrico es una función de la forma:
\begin{equation*}
    f(t) = \sum_{n=M}^N c_n e^{int} = \frac{a_0}{2} + \sum_{n=1}^{\max(N,M)} (a_n \cos(nt) + b_n \sin(nt))
\end{equation*}
Denotamos al conjunto de polinomios trigonométricos como $\mathcal{P}(\mathbb{T})$.
\end{definicion}

\begin{observacion}{Recordatorio de Topología y Densidad}{recordatorio_densidad}
Supongamos un espacio topológico $X$, con $A, B \subseteq X$.
\begin{itemize}
    \item $A$ es denso $\iff \bar{A} = X$.
    \item Si $X$ es normado, $A$ es denso $\iff$ Fijado $x \in X, \forall \epsilon > 0 \ \exists y \in A \text{ tq } \|x - y\| < \epsilon$.
    \item Si $A$ es denso en $B$ y $B$ es denso en $C \implies A$ es denso en $C$.
    \item Si $A$ es denso en $C$ y $A \subseteq B \subseteq C \implies B$ es denso en $C$.
\end{itemize}

Si $C(K) = \{f: K \to \mathbb{C} \mid \text{continua}\}$, un subconjunto $A \subseteq C(K)$ cumple:
\begin{enumerate}
    \item $f_1, f_2 \in A \implies f_1 + f_2 \in A$
    \item $f_1, f_2 \in A \implies f_1 \cdot f_2 \in A$
\end{enumerate}
Entonces $A$ forma una \textbf{subálgebra}. Si además se cumple que:
\begin{enumerate}
    \setcounter{enumi}{2}
    \item $f \in A \implies \bar{f} \in A$
\end{enumerate}
Entonces $A$ es denso en $C(K)$.
\end{observacion}

\begin{proposicion}{Densidad en $C(\mathbb{T})$}{densidad_ct}
En nuestro caso $C(K) = C(\mathbb{T})$.

Si $f, g$ son polinomios trigonométricos:
\begin{align*}
    f(t) &= \sum_{n=-N_1}^{N_1} c_n e^{int} \\
    g(t) &= \sum_{n=-N_2}^{N_2} d_n e^{int}
\end{align*}
Se cumplen de manera directa las siguientes propiedades:
\begin{enumerate}
    \item $f + g \in \mathcal{P}(\mathbb{T})$
    \item $f \cdot g \in \mathcal{P}(\mathbb{T})$
    \item $\bar{f} = \sum_{n=-N_1}^{N_1} \overline{c_n e^{int}} = \sum_{n=-N_1}^{N_1} \overline{c_n} \cdot e^{-int} \in \mathcal{P}(\mathbb{T})$
\end{enumerate}
Por tanto, $\mathcal{P}(\mathbb{T})$ es denso en $C(\mathbb{T})$.
\end{proposicion}

\begin{nota}{Interpretación del Teorema de Densidad}{interpretacion_prop}
El resultado de la proposición anterior quiere decir:

Para toda $f \in C(\mathbb{T})$ y todo $\epsilon > 0$, existe un $P \in \mathcal{P}(\mathbb{T})$ (polinomio trigonométrico) tal que:
\begin{equation*}
    \sup_{t \in [-\pi, \pi]} |P(t) - f(t)| \le \epsilon
\end{equation*}
Es decir, $\|f - P\|_\infty \le \epsilon$.
\end{nota}
\begin{definicion}{El espacio $L_p(\mathbb{T})$}{espacio_lp}
Para $1 \le p < \infty$, decimos que una función $2\pi$-periódica $f:\mathbb{R} \to \mathbb{C}$ pertenece al espacio $L_p(\mathbb{T})$ si es medible en el sentido de Lebesgue y la integral de su valor absoluto a la potencia $p$ es finita. Formalmente:
\begin{equation*}
    f \in L_p(\mathbb{T}) \iff \text{medible Lebesgue y } \underbrace{\left( \frac{1}{2\pi} \int_{-\pi}^{\pi} |f(t)|^p dt \right)^{1/p}}_{\|f\|_p} < \infty
\end{equation*}
\end{definicion}

\begin{nota}{Propiedad en casi todas partes (c.p.p.)}{cpp}
Decimos que una propiedad matemática se cumple \textbf{casi por todas partes} (c.p.p. o \textit{a.e.} en inglés) si el conjunto de puntos donde no se verifica es tan insignificante que tiene medida de Lebesgue igual a cero.
\end{nota}

\begin{observacion}{Clases de equivalencia}{clases_eq}
En la práctica funcional, no distinguimos entre funciones que difieren en un conjunto de medida cero. Trabajamos sistemáticamente con clases de equivalencia, asumiendo que $f = g$ si coinciden casi por todas partes ($f \sim g$).
\end{observacion}

\begin{teorema}{Desigualdad de Hölder}{holder}
Imaginemos dos funciones $f \in L_r(\mathbb{T})$ y $g \in L_{r'}(\mathbb{T})$, donde los exponentes son conjugados (cumplen la relación $\frac{1}{r} + \frac{1}{r'} = 1$). Entonces, su producto $f \cdot g$ siempre será integrable ($L_1$), y su norma quedará acotada por el producto de sus normas respectivas:
\begin{equation*}
    \frac{1}{2\pi} \int_{-\pi}^{\pi} |f(t)g(t)| dt \le \left( \frac{1}{2\pi} \int_{-\pi}^{\pi} |f(t)|^r dt \right)^{1/r} \left( \frac{1}{2\pi} \int_{-\pi}^{\pi} |g(t)|^{r'} dt \right)^{1/r'}
\end{equation*}
De forma más compacta, esto se escribe como $\|fg\|_1 \le \|f\|_r \cdot \|g\|_{r'}$.
\end{teorema}

\begin{proposicion}{Inclusión de espacios $L_p$ y densidad}{incl_dens_lp}
Se cumplen las siguientes propiedades fundamentales para los espacios $L_p(\mathbb{T})$:
\begin{enumerate}
    \item \textbf{Inclusión de espacios:} Si un exponente es estrictamente menor que otro ($p < q$), el espacio $L_q$ está totalmente contenido en $L_p$. Es decir, si $\|f\|_q < \infty$, entonces $\|f\|_p < \infty$.
    
    \item \textbf{Caso del límite infinito:} Si asumimos que la función es continua en el toro ($f \in C(\mathbb{T})$) y por tanto acotada, su norma suprema domina a todas las demás normas $p$:
    \begin{equation*}
        \left( \frac{1}{2\pi} \int_{-\pi}^{\pi} |f(t)|^p dt \right)^{1/p} \le \left( \|f\|_\infty^p \frac{1}{2\pi} \int_{-\pi}^{\pi} dt \right)^{1/p} = \|f\|_\infty
    \end{equation*}
    
    \item \textbf{Densidad:} Las funciones continuas $C(\mathbb{T})$ forman un conjunto denso en $L_1(\mathbb{T})$ (y, en general, en cualquier $L_p(\mathbb{T})$). Dicho de otra manera: dada una $f \in L_1$ y cualquier margen de error $\varepsilon > 0$, siempre podremos encontrar una función continua $g \in C(\mathbb{T})$ tal que $\|f - g\|_1 < \varepsilon$. 
    
    Como conclusión general de este apartado, la jerarquía de densidades es $\mathcal{P}(\mathbb{T}) \subseteq C(\mathbb{T}) \subseteq L_1(\mathbb{T})$.
\end{enumerate}
\end{proposicion}

\begin{proof}
\textbf{Demostración de (1):} Expresamos la integral de $|f|^p$ como el producto $|f|^p \cdot 1$. Aplicamos la desigualdad de Hölder con el exponente $r = q/p$ y su respectivo conjugado $r'$:
\begin{align*}
    \frac{1}{2\pi} \int_{-\pi}^{\pi} |f(t)|^p dt &= \frac{1}{2\pi} \int_{-\pi}^{\pi} |f(t)|^p \cdot 1 dt \\
    &\le \left( \frac{1}{2\pi} \int_{-\pi}^{\pi} \left(|f(t)|^p\right)^r dt \right)^{1/r} \cdot \left( \frac{1}{2\pi} \int_{-\pi}^{\pi} 1^{r'} dt \right)^{1/r'} \\
    &= \left( \frac{1}{2\pi} \int_{-\pi}^{\pi} |f(t)|^q dt \right)^{p/q} \cdot 1
\end{align*}
Elevando ambos lados a la potencia $1/p$, llegamos a la acotación directa $\|f\|_p \le \|f\|_q$.

\medskip
\textbf{Demostración de (3) (intuitiva en 3 pasos):}
\begin{itemize}
    \item \textbf{Caso 1:} Si $f$ es una función indicadora estricta $f = \chi_{[a,b]}$, construimos un trapecio continuo $g_\delta$ que valga $1$ en $[a,b]$ y caiga a $0$ de manera lineal a lo largo de un pequeño margen $\delta$. El área de discrepancia será $\|f - g_\delta\|_1 \le \frac{\delta}{\pi}$. Nos basta con tomar $\delta = \pi \varepsilon$.
    \item \textbf{Caso 2:} Si $f$ es una función simple (combinación lineal finita de indicadores sobre intervalos disjuntos), aplicamos la técnica del trapecio a cada escalón.
    \item \textbf{Caso 3:} Si $f \in L_1$ es cualquier función integrable, siempre se puede aproximar primero por una función simple $s$ tal que $\|f - s\|_1 < \varepsilon/2$. Por el Caso 2, aproximamos esa $s$ con una continua $g$ a distancia $\varepsilon/2$. Por la desigualdad triangular, el error total es menor a $\varepsilon$.
\end{itemize}
\end{proof}

\begin{proposicion}{Coeficientes de un polinomio trigonométrico}{coef_pol_trig}
Dado un polinomio trigonométrico $P(t) = \sum_{k=-M}^M c_k e^{ikt}$, sus coeficientes de Fourier $\hat{P}(n)$ coinciden de forma exacta con los propios coeficientes $c_n$.

\textit{Demostración:} Al tratarse de una suma finita, la linealidad de la integral nos permite intercambiar sin problemas el sumatorio y la integración:
\begin{equation*}
    \hat{P}(n) = \frac{1}{2\pi} \int_{-\pi}^{\pi} \left( \sum_{k=-M}^M c_k e^{ikt} \right) e^{-int} dt = \sum_{k=-M}^M c_k \left( \frac{1}{2\pi} \int_{-\pi}^{\pi} e^{i(k-n)t} dt \right)
\end{equation*}
Por la ortogonalidad natural del sistema, la integral vale $1$ si $k=n$ y $0$ en el resto de casos, colapsando toda la suma a:
\begin{equation*}
    \hat{P}(n) = \begin{cases} c_n & \text{si } -M \le n \le M \\ 0 & \text{en cualquier otro caso} \end{cases}
\end{equation*}
\end{proposicion}

\begin{observacion}{Recuperación mutua de coeficientes reales y complejos}{rel_coef_real_comp}
Si tienes entre manos los coeficientes complejos $(\hat{f}(n))_{n \in \mathbb{Z}}$, puedes reconstruir inmediatamente los coeficientes de la serie real clásica ($\frac{a_0}{2} + \sum [a_n \cos(nt) + b_n \sin(nt)]$) mediante estas tres expresiones:
\begin{align*}
    a_0 &= 2\hat{f}(0) \\
    a_n &= \hat{f}(n) + \hat{f}(-n) \\
    b_n &= i(\hat{f}(n) - \hat{f}(-n))
\end{align*}
\end{observacion}

\begin{nota}{Ejemplo práctico: Coeficientes de la función signo}{ejemplo_signo}
Vamos a calcular explícitamente los coeficientes de Fourier $\hat{f}(n)$ para la función signo en el toro, definida como $f(t) = 1$ para $t \in [0, \pi)$ y $f(t) = -1$ para $t \in [-\pi, 0)$.

Separamos la integral de evaluación en los dos hemiciclos:
\begin{equation*}
    \hat{f}(n) = \frac{1}{2\pi} \int_{-\pi}^{0} (-1) e^{-int} dt + \frac{1}{2\pi} \int_{0}^{\pi} (1) e^{-int} dt
\end{equation*}

Evaluando el término central nulo ($n = 0$), las áreas se anulan perfectamente:
\begin{equation*}
    \hat{f}(0) = -\frac{\pi}{2\pi} + \frac{\pi}{2\pi} = 0
\end{equation*}

Para cualquier frecuencia no nula ($n \ne 0$), integramos directamente la exponencial:
\begin{align*}
    \hat{f}(n) &= -\frac{1}{2\pi} \left[ \frac{e^{-int}}{-in} \right]_{-\pi}^{0} + \frac{1}{2\pi} \left[ \frac{e^{-int}}{-in} \right]_{0}^{\pi} \\
    &= \frac{1}{2\pi i n} \left( 1 - e^{in\pi} \right) - \frac{1}{2\pi i n} \left( e^{-in\pi} - 1 \right)
\end{align*}

Sabiendo que la identidad de Euler nos dicta que $e^{-in\pi} = e^{in\pi} = (-1)^n$, podemos agrupar ambos corchetes:
\begin{equation*}
    \hat{f}(n) = \frac{1}{\pi i n} \left[ 1 - (-1)^n \right] = \begin{cases} 
        0 & \text{si } n \text{ es par } (n=2k) \\ 
        -\frac{2i}{\pi n} & \text{si } n \text{ es impar } (n=2k+1) 
    \end{cases}
\end{equation*}
Esto nos genera una serie puramente impar, expresada como:
\begin{equation*}
    S(f)(t) = \sum_{k \in \mathbb{Z}} \frac{-2i}{(2k+1)\pi} e^{i(2k+1)t}
\end{equation*}
\end{nota}

\begin{definicion}{El Coeficiente y la Serie de Fourier}{def_fourier_formal}
Para toda función integrable $f \in L_1(\mathbb{T})$, definimos el \textbf{coeficiente de Fourier $n$-ésimo} (con $n \in \mathbb{Z}$) de la siguiente manera:
\begin{equation*}
    \hat{f}(n) = \frac{1}{2\pi} \int_{-\pi}^{\pi} f(t) e^{-int} dt
\end{equation*}
Esta operación está siempre bien definida porque la acotación modular $|f(t)e^{-int}| = |f(t)|$ garantiza que el integrando permanece dentro de las fronteras de $L_1$, produciendo que $\hat{f}(n) \in \mathbb{C}$ exista de manera finita.

Con la colección de estos coeficientes generamos la \textbf{Serie de Fourier} asociada a $f$:
\begin{equation*}
    f(t) \sim \sum_{n=-\infty}^{\infty} \hat{f}(n) e^{int}
\end{equation*}
Para poder manejar y analizar esta serie infinita definimos la \textbf{suma parcial n-ésima}, $S(f,n)(t) = \sum_{k=-n}^{n} \hat{f}(k) e^{ikt}$, que resulta ser un polinomio trigonométrico que mapea $[-\pi,\pi] \to \mathbb{C}$.
\end{definicion}

\begin{observacion}{Dudas sobre la convergencia}{conv_uniforme_fourier}
\textbf{Pregunta clave:} ¿Podemos asumir que la suma parcial $S(f,n)$ acabará convergiendo (y además uniformemente) de regreso hacia nuestra $f$ original? 

La respuesta es un rotundo \textbf{no}. La razón profunda es topológica: si el límite uniforme de los polinomios trigonométricos $S(f,n)$ existiese siempre, forzosamente daría a luz una función continua. Sin embargo, el espacio $L_1(\mathbb{T})$ está completamente repleto de funciones discontinuas, por lo que la convergencia uniforme general es imposible de garantizar sin hipótesis de regularidad añadidas.
\end{observacion}

\begin{nota}{Ortogonalidad de la base exponencial}{ejemplo_base_exp}
Para ver cómo funciona la maquinaria sobre sus propios cimientos, analicemos la función base $\phi_m(t) = e^{imt}$. Su coeficiente de Fourier $\hat{\phi}_m(n)$ evaluará la presencia de la frecuencia $n$ en la frecuencia $m$:
\begin{equation*}
    \hat{\phi}_m(n) = \frac{1}{2\pi} \int_{-\pi}^{\pi} e^{imt} e^{-int} dt = \frac{1}{2\pi} \int_{-\pi}^{\pi} e^{i(m-n)t} dt
\end{equation*}

\begin{itemize}
    \item \textbf{Si $m = n$:} la integral colapsa a la unidad básica $\frac{1}{2\pi} \int_{-\pi}^{\pi} e^0 dt = 1$.
    \item \textbf{Si $m \ne n$:} llamando a su diferencia no nula $k = m-n$, aplicamos la primitiva:
    \begin{equation*}
        \hat{\phi}_m(n) = \frac{1}{2\pi i k} \left[ e^{ikt} \right]_{-\pi}^{\pi} = \frac{1}{2\pi i k} \left( e^{ik\pi} - e^{-ik\pi} \right) = 0
    \end{equation*}
\end{itemize}

El resultado es la clásica Delta de Kronecker, demostrando que las funciones armónicas filtran sus frecuencias de manera ortonormal:
\begin{equation*}
    \hat{\phi}_m(n) = \begin{cases} 
        1 & \text{si } n=m \\ 
        0 & \text{si } n \ne m 
    \end{cases}
\end{equation*}
\end{nota}